\section{Problem Formulation}\label{sec:problem}
The long-only mean-variance portfolio solves
\[
\min_{w}\; w^\top\Sigma\, w - \rho\,\mu^\top w
\quad\text{subject to}\quad
Bw = c,\quad w\geq 0~,
\]
where $\Sigma$ is the $n\times n$ covariance matrix of returns, $\mu$ is the $n\times 1$ vector of expected returns, $\rho\geq 0$ controls the risk-return trade-off (setting $\rho=0$ recovers the pure long-only minimum variance portfolio), and $B \in \mathbb{R}^{p\times n}$ of full row rank, $c \in \mathbb{R}^p$ carry $p \ll n$ linear balance conditions. The budget constraint $\mathbf{1}^\top w = 1$ (full investment) is the case $p=1$, $B = \mathbf{1}^\top$, $c = 1$; it is the constraint used in every experiment reported here, and the reader may substitute it throughout. The general $B$ additionally accommodates group budgets (sector sleeves summing to prescribed weights), factor neutralities $\beta^\top w = 0$, or a target portfolio beta at the same cost structure. This is the equality-augmented non-negative quadratic of the companion paper~\cite{schmelzer2026nncg}, whose treatment the present formulation follows.
Since both $\Sigma$ and
$\mu$ are unknown, in practice they must be replaced by estimates. Let $X\in\mathbb{R}^{T\times n}$ denote the demeaned return matrix for $n$ assets over the previous $T$ trading periods. The feasible problem formulation then is
\[
\min_{w}\; w^\top\hat\Sigma\, w - \rho\,\hat\mu^\top w
\quad\text{subject to}\quad
Bw = c,\quad w\geq 0~.
\]
The classical choice for estimating $\Sigma$ is
the sample covariance matrix $\hat\Sigma = X^\top X/T$, for which $w^\top\hat\Sigma\, w = \norm{Xw}^2/T$, and the problem becomes
\[
\min_{w}\;\frac{\norm{Xw}^2}{T} - \rho\,\hat\mu^\top w
\quad\text{subject to}\quad
Bw = c,\quad w\geq 0~.
\]
This reframing is key: the variance term is now a normalised squared norm of $Xw$, and the gradient $w\mapsto X^\top(Xw)$ can be evaluated with two matrix-vector products, without ever forming $X^\top X$.

\begin{table}[ht]
\centering
\begin{tabular}{lll}
\toprule
Parameter & Type & Description \\
\midrule
$X$        & $\mathbb{R}^{T \times n}$ & Demeaned return matrix \\
$\rho$     & $\mathbb{R}_{\geq 0}$     & Return-tilt weight (default $0$) \\
$\mu$      & $\mathbb{R}^n$            & Expected returns (default $\mathbf{0}$) \\
$B$        & $\mathbb{R}^{p \times n}$ & Balance constraints, full row rank (default $\mathbf{1}^\top$) \\
$c$        & $\mathbb{R}^{p}$          & Balance targets (default $1$) \\
\bottomrule
\end{tabular}
\caption{Parameters of the \texttt{Problem} specification.}
\label{tab:params}
\end{table}

\paragraph{KKT structure}
Up to the factor $\tfrac{1}{2}$ in the quadratic term, this is the equality-augmented
non-negative quadratic of the companion paper~\cite{schmelzer2026nncg} with
$A = \tfrac{2}{T}X^\top X$ and $b = \rho\hat\mu$; its KKT, complementarity, and
$P$-matrix pivot structure are developed there in full and need only be read in
portfolio terms. Introducing multipliers $\lambda \in \mathbb{R}^p$ for the balance
constraints and $\nu \geq \mathbf{0}$ for non-negativity, stationarity reads
\[
\tfrac{2}{T}\, X^\top X w = B^\top\lambda+\nu+\rho\hat\mu\;,
\]
and complementary slackness ($\nu_i w_i = 0$ for all $i$) partitions assets
into two types. \emph{Active} assets ($w_i>0$) have $\nu_i=0$; stationarity then gives
\[
\tfrac{2}{T}\,(X^\top Xw)_i = (B^\top\lambda)_i+\rho\hat\mu_i\;,
\]
pinning each active asset's marginal variance, adjusted for its expected return, to the
row space of $B$ (in the budget case, to the single shadow price $\lambda$ of full
investment). \emph{Excluded} assets ($w_i=0$) satisfy
\[
\nu_i = \tfrac{2}{T}\,(X^\top Xw)_i - (B^\top\lambda)_i - \rho\hat\mu_i \geq 0~.
\]
Their marginal variance exceeds this threshold, so holding them at a positive weight would
increase variance net of return. Once the estimator is strictly positive definite
(guaranteed by shrinkage with $\alpha > 0$, Section~\ref{sec:precond}), the objective is
strictly convex and these conditions are necessary and sufficient for the unique global
optimum \cite[Section~2]{schmelzer2026nncg}. The primal-dual loop in
Section~\ref{sec:primaldual} is a direct
implementation of these conditions: it terminates precisely when primal and dual feasibility
both hold, which (together with stationarity from each inner solve) is sufficient for
global optimality.
